\chapter{Simulations and results}

This chapter compares the algorithms implemented in this work against the
other Shortest Common Superstring heuristics and approximation algorithms
already present in the repository\footnote{\url{\repoUrl}}: the classical
\textsc{Greedy} (maximum overlap) heuristic with its \textsc{MGreedy} and
\textsc{TGreedy} variants, Group
Merge, the cycle-cover algorithms of Teng and
Yao~\cite{teng1997approximating} and of Paluch, Elbassioni and van
Zuylen~\cite{paluch2012simpler}, and the $\frac{38}{63}$-approximation
directed Max-TSP algorithm of Kosaraju, Park and Stein~\cite{Kosaraju1994}
used as an overlap approximation subroutine.

Four of the nine implementations benchmarked here were contributed as part
of this work\footnote{\url{\repoUrl/pull/61}}: the two
Breslauer--Jiang--Jiang algorithms of
Chapter~\ref{chap:breslauer}, the Kosaraju--Park--Stein Max-TSP
approximation of Chapter~\ref{chap:kosaraju} that the second of them uses
as a subroutine, and the path-coloring algorithm of
Chapter~\ref{chap:path_coloring} that Version~2 of Kosaraju--Park--Stein
requires in turn. Table~\ref{tab:impl_locations} lists where each of them
is implemented, with the contributions of this work set in bold.

\begin{table}[htbp]
    \centering
    \caption{Where each benchmarked algorithm is implemented, as a path
    relative to the root of the repository. The implementations contributed
    by this work are set in \textbf{bold}.}
    \label{tab:impl_locations}
    \footnotesize
    \setlength{\tabcolsep}{5pt}
    \renewcommand{\arraystretch}{1.2}
    \begin{tabular}{@{}>{\raggedright\arraybackslash}p{135pt}l@{}}
        \toprule
        \textbf{Algorithm} & \textbf{File} \\
        \midrule
        \textsc{Greedy}, \textsc{MGreedy}, \textsc{TGreedy}, Group Merge, brute-force optimum & \texttt{\detokenize{shortest_common_superstring/shortest_common_superstring.py}} \\
        Teng--Yao & \texttt{\detokenize{shortest_common_superstring/teng_yao.py}} \\
        Paluch--Elbassioni--van~Zuylen & \texttt{\detokenize{shortest_common_superstring/paluch_elbassoni_zuylen.py}} \\
        \addlinespace[3pt]
        \textbf{Kosaraju--Park--Stein ($\frac{38}{63}$)} & \texttt{\detokenize{shortest_common_superstring/kosaraju.py}} \\
        \textbf{Path-coloring lemma} & \texttt{\detokenize{shortest_common_superstring/path_coloring.py}} \\
        \textbf{Breslauer--Jiang--Jiang $2\frac{2}{3}$ and $2\frac{25}{42}$} & \texttt{\detokenize{shortest_common_superstring/breslauer.py}} \\
        \bottomrule
    \end{tabular}
\end{table}

\section{Method}

No standard public benchmark suite exists for the Shortest Common
Superstring problem, unlike, e.g., TSPLIB for the Traveling Salesman
Problem. Three sources of test instances were therefore used.

\begin{itemize}[itemsep=3pt,topsep=4pt]
\item \textbf{Random strings.} Independent words over an alphabet of
      controlled size, with lengths drawn uniformly from $8$ to $15$.
\item \textbf{Natural text.} Fragments of length $25$ cut at random
      positions from \emph{Pan Tadeusz}, lowercased and stripped of
      punctuation.
\item \textbf{DNA.} Fragments of length $40$ cut at random positions from
      the genome of bacteriophage phiX174~\cite{sanger1977phix174}, NCBI
      accession NC\_001422.1~\cite{ncbi_phix174}.
\end{itemize}

The last two follow the random-fragment methodology used in the DNA
assembly literature: a single longer source string is cut into overlapping
pieces at random positions, rather than generating fully independent random
words.

Because the true optimum is computable by brute force only for very small
instances, each result is reported relative to the shortest superstring
found by any of the algorithms on the same instance. A value of $1.0000$
therefore means the algorithm was never beaten on any instance of that
experiment. Section~\ref{sec:bench_exact} covers instances small enough for
the brute-force solver and compares against the true optimum instead.

Every reported value is a mean over independent repetitions of the same
experiment. Running times were measured on a single core with the overlap
matrix computed inside the timed region, so they include the cost of
building the overlap graph.

\section{Overview}

Table~\ref{tab:bench_overview} pools the four comparative experiments of
the following sections into one ranking, over all instance sizes, alphabets
and sources at once. The experiment on tiny instances is excluded from it,
because there the reference is the true optimum rather than the best result
found.

\begin{table}[htbp]
    \centering
    \caption{Overview of all benchmarked algorithms over the four comparative experiments: superstring length relative to the best result found on the same instance, and mean running time.}
    \label{tab:bench_overview}
    \small
    \setlength{\tabcolsep}{7pt}
    \renewcommand{\arraystretch}{1.1}
    \begin{tabular}{@{}lrrr@{}}
        \toprule
        \textbf{Algorithm} & \multicolumn{2}{c}{\textbf{Length / best found}} & \textbf{Mean time} \\
        \cmidrule(lr){2-3}
         & mean & worst & (ms) \\
        \midrule
        \textsc{Greedy} (max overlap) & 1.0052 & 1.0364 & 28.42 \\
        \textsc{MGreedy} & 1.0115 & 1.0992 & 28.17 \\
        \textsc{TGreedy} & 1.0071 & 1.0577 & 29.77 \\
        Group Merge & 1.0093 & 1.0534 & 6.60 \\
        Teng--Yao & 1.0191 & 1.1264 & 5.58 \\
        Paluch--Elbassioni--van~Zuylen & 1.0018 & 1.0221 & 3.42 \\
        Kosaraju--Park--Stein ($\frac{38}{63}$) & 1.0046 & 1.1074 & 17.31 \\
        Breslauer--Jiang--Jiang $2\frac{2}{3}$ & 1.0794 & 1.3158 & 2.99 \\
        Breslauer--Jiang--Jiang $2\frac{25}{42}$ & 1.0772 & 1.3158 & 3.03 \\
        \bottomrule
    \end{tabular}
\end{table}


\section{Scaling with instance size and alphabet}

Figure~\ref{fig:bench_ratio_vs_n} varies the number of input strings at a
fixed alphabet $\{a,b,c,d\}$, and Figure~\ref{fig:bench_ratio_vs_alphabet}
varies the alphabet size at a fixed $n=20$.
Figure~\ref{fig:bench_runtime_vs_n} reports the running times measured on
the instances of Figure~\ref{fig:bench_ratio_vs_n}.

\begin{figure}[tbp]
    \centering
    \includegraphics[width=0.92\textwidth]{chart_ratio_vs_n}
    \caption{Mean superstring length relative to the best result found, as a
    function of the number of input strings $n$ (random strings, alphabet
    $\{a,b,c,d\}$, word lengths $8$--$15$, 8 repeats per point).}
    \label{fig:bench_ratio_vs_n}
\end{figure}

\begin{figure}[tbp]
    \centering
    \includegraphics[width=0.92\textwidth]{chart_ratio_vs_alphabet}
    \caption{Mean superstring length relative to the best result found, as a
    function of alphabet size (random strings, $n=20$, word lengths
    $8$--$15$, 8 repeats per point).}
    \label{fig:bench_ratio_vs_alphabet}
\end{figure}

\begin{figure}[tbp]
    \centering
    \includegraphics[width=0.92\textwidth]{chart_runtime_vs_n}
    \caption{Mean running time (log scale) as a function of the number of
    input strings $n$, same instances as Figure~\ref{fig:bench_ratio_vs_n}.}
    \label{fig:bench_runtime_vs_n}
\end{figure}

\section{Real-world data: natural text and DNA}

Figure~\ref{fig:bench_source_comparison} puts the three instance sources
side by side at the same relative scale.
Table~\ref{tab:bench_dna} reports absolute lengths on the DNA instances
rather than ratios, together with the compression achieved relative to the
plain concatenation of all reads.

\begin{figure}[tbp]
    \centering
    \includegraphics[width=0.92\textwidth]{chart_source_comparison}
    \caption{Mean superstring length relative to the best result found,
    grouped by instance source: random strings, fragments of a
    natural-language text (\emph{Pan Tadeusz}), and fragments of a real
    genome (bacteriophage phiX174).}
    \label{fig:bench_source_comparison}
\end{figure}

\begin{table}[htbp]
    \centering
    \caption{Mean superstring length on real DNA fragments (bacteriophage phiX174, NCBI NC\_001422.1~\cite{ncbi_phix174}), averaged over $k \in \{8,15,25,40,60\}$ reads of length 40 at 5 repeats each. The last column is the length divided by the length of the plain concatenation of all reads, i.e.\ the compression actually achieved.}
    \label{tab:bench_dna}
    \small
    \setlength{\tabcolsep}{7pt}
    \renewcommand{\arraystretch}{1.1}
    \begin{tabular}{@{}lrr@{}}
        \toprule
        \textbf{Algorithm} & \textbf{Mean superstring length} & \textbf{vs.\ concatenation} \\
        \midrule
        \textsc{Greedy} (max overlap) & 986.1 & 0.833 \\
        \textsc{MGreedy} & 988.0 & 0.834 \\
        \textsc{TGreedy} & 986.8 & 0.833 \\
        Group Merge & 991.4 & 0.837 \\
        Teng--Yao & 998.5 & 0.843 \\
        Paluch--Elbassioni--van~Zuylen & 985.4 & 0.832 \\
        Kosaraju--Park--Stein ($\frac{38}{63}$) & 994.4 & 0.840 \\
        Breslauer--Jiang--Jiang $2\frac{2}{3}$ & 1044.9 & 0.883 \\
        Breslauer--Jiang--Jiang $2\frac{25}{42}$ & 1044.6 & 0.882 \\
        \addlinespace[2pt]
        \textit{Plain concatenation (no overlap)} & 1184.0 & 1.000 \\
        \bottomrule
    \end{tabular}
\end{table}


\section{Comparison against the true optimum}
\label{sec:bench_exact}

The brute-force solver implemented alongside these algorithms enumerates
candidate superstrings by increasing length, which is tractable only on
very small inputs. Table~\ref{tab:bench_exact_small} therefore uses 20
independent instances of three or four words of length three or four over a
binary alphabet, on which the solver still terminates quickly, and reports
every algorithm against the true optimum rather than against the best
result found.

\begin{table}[htbp]
    \centering
    \caption{Comparison against the true optimum, computed by brute-force enumeration. A single instance consists of $n \in \{3,4\}$ words of length $3$ or $4$ over the alphabet $\{a,b\}$, and the table averages over 20 such instances.}
    \label{tab:bench_exact_small}
    \small
    \setlength{\tabcolsep}{7pt}
    \renewcommand{\arraystretch}{1.1}
    \begin{tabular}{@{}lrrr@{}}
        \toprule
        \textbf{Algorithm} & \multicolumn{2}{c}{\textbf{Length / OPT}} & \textbf{Optimal} \\
        \cmidrule(lr){2-3}
         & mean & worst & instances \\
        \midrule
        \textsc{Greedy} (max overlap) & 1.1541 & 1.2500 & 0/20 \\
        \textsc{MGreedy} & 1.1091 & 1.5000 & 11/20 \\
        \textsc{TGreedy} & 1.1541 & 1.2500 & 0/20 \\
        Group Merge & 1.0000 & 1.0000 & 20/20 \\
        Teng--Yao & 1.0872 & 1.3333 & 11/20 \\
        Paluch--Elbassioni--van~Zuylen & 1.0083 & 1.1667 & 19/20 \\
        Kosaraju--Park--Stein ($\frac{38}{63}$) & 1.0674 & 1.3750 & 14/20 \\
        Breslauer--Jiang--Jiang $2\frac{2}{3}$ & 1.1488 & 1.4286 & 9/20 \\
        Breslauer--Jiang--Jiang $2\frac{25}{42}$ & 1.1488 & 1.4286 & 9/20 \\
        \bottomrule
    \end{tabular}
\end{table}


\section{The cost of a single cycle-breaking version}
\label{sec:bench_kosaraju_versions}

The Kosaraju--Park--Stein algorithm of Chapter~\ref{chap:kosaraju} computes
one maximum weight cycle cover and then breaks it into paths in three
independent ways. Version~1 discards the lightest edge of every cycle
(Section~\ref{sec:kps_v1}), Version~2 contracts the $2$-cycles and applies
the path-coloring lemma of Chapter~\ref{chap:path_coloring}
(Section~\ref{sec:kps_v2}), and Version~3 cuts long cycles into pieces and
recombines them through a maximum weight matching
(Section~\ref{sec:kps_v3}). The algorithm returns the heaviest of the three
tours, and the $\frac{38}{63}$ guarantee of
Theorem~\ref{thm:final_approximation} is a statement about that maximum.

To measure what each version contributes on its own, the same algorithm was
run four times per instance with the version set restricted to $\{1\}$,
$\{2\}$, $\{3\}$ and $\{1,2,3\}$. All four start from the identical cycle
cover, so every difference comes from the cycle-breaking step alone.
Lengths in Table~\ref{tab:bench_kosaraju_versions} and
Figure~\ref{fig:bench_kosaraju_versions} are relative to the best of these
four variants, not to the full set of algorithms benchmarked earlier.

\begin{table}[htbp]
    \centering
    \caption{Effect of restricting the Kosaraju--Park--Stein algorithm to a single cycle-breaking version instead of taking the heaviest of the three. Lengths are relative to the best of these four variants on the same instance. The random columns cover $n \in \{5,\dots,50\}$ (8 repeats per point), the DNA column covers phiX174 reads with $k \in \{8,15,25,40\}$ (5 repeats per point).}
    \label{tab:bench_kosaraju_versions}
    \small
    \setlength{\tabcolsep}{6pt}
    \renewcommand{\arraystretch}{1.1}
    \begin{tabular}{@{}lrrrrr@{}}
        \toprule
        \textbf{Variant} & \multicolumn{2}{c}{\textbf{Random strings}} & \textbf{DNA} & \textbf{Shortest} & \textbf{Mean time} \\
        \cmidrule(lr){2-3} \cmidrule(lr){4-4}
         & mean & worst & mean & found & (ms) \\
        \midrule
        Version~1 only & 1.0006 & 1.0128 & 1.0004 & 76/84 & 1.38 \\
        Version~2 only & 1.0365 & 1.0773 & 1.0284 & 7/84 & 4.94 \\
        Version~3 only & 1.0043 & 1.0213 & 1.0009 & 48/84 & 11.25 \\
        \addlinespace[2pt]
        Best of 1--3 (full algorithm) & 1.0000 & 1.0000 & 1.0000 & 84/84 & 14.81 \\
        \bottomrule
    \end{tabular}
\end{table}


\begin{figure}[htbp]
    \centering
    \includegraphics[width=\textwidth]{chart_kosaraju_versions}
    \caption{Quality and running time of the Kosaraju--Park--Stein algorithm
    when restricted to a single cycle-breaking version, as a function of the
    number of input strings $n$ (random strings, alphabet $\{a,b,c,d\}$,
    word lengths $8$--$15$, 8 repeats per point). Lengths are relative to
    the best of the four variants on the same instance.}
    \label{fig:bench_kosaraju_versions}
\end{figure}

\section{Large instances}
\label{sec:bench_large}

The experiments above stay below $n = 60$, where every algorithm finishes in
milliseconds and the running times compare implementation constants rather
than asymptotic behaviour. This section repeats the comparison on instances
large enough for the cubic step to dominate. Fragments were cut from
\emph{Pan Tadeusz} at $n \in \{200, 400, 800, 1600\}$ and from the genome of
phage lambda~\cite{sanger1982lambda}, NCBI accession
NC\_001416.1~\cite{ncbi_lambda}, at $n \in \{200, 400, 800\}$ with reads of
length 100. Each point is a single instance handed to all nine
algorithms, and each run was given a 90-minute limit, which none of them
reached.

\begin{figure}[tbp]
    \centering
    \includegraphics[width=\textwidth]{chart_large_scale}
    \caption{Running time on large instances, both axes logarithmic, so the
    slope of each line is the empirical exponent. Fitted on the text
    instances, the three greedy variants and Group Merge follow slopes of
    3.0 and 3.1, Teng--Yao,
    Paluch--Elbassioni--van~Zuylen and the two Breslauer--Jiang--Jiang
    algorithms fall between 2.0 and
    2.1, and Kosaraju--Park--Stein sits between the two
    groups at 2.8. \textsc{Greedy}, \textsc{MGreedy} and
    \textsc{TGreedy} differ by at most 6 percent at any
    point and are drawn on top of one another.}
    \label{fig:bench_large_scale}
\end{figure}

\begin{table}[htbp]
    \centering
    \caption{The largest instance measured for each source: superstring length relative to the shortest found on that instance, and running time. Every algorithm was given the same instance.}
    \label{tab:bench_large_scale}
    \small
    \setlength{\tabcolsep}{6pt}
    \renewcommand{\arraystretch}{1.1}
    \begin{tabular}{@{}lrrrr@{}}
        \toprule
        \textbf{Algorithm} & \multicolumn{2}{c}{\textbf{Text, $n=1600$}} & \multicolumn{2}{c}{\textbf{DNA, $n=800$}} \\
        \cmidrule(lr){2-3} \cmidrule(lr){4-5}
          & length / best & time (s) & length / best & time (s) \\
        \midrule
        \textsc{Greedy} (max overlap) & 1.0000 & 4152.5 & 1.0000 & 1421.3 \\
        \textsc{MGreedy} & 1.0004 & 4091.3 & 1.0003 & 1420.8 \\
        \textsc{TGreedy} & 1.0002 & 4090.5 & 1.0001 & 1422.1 \\
        Group Merge & 1.0054 & 495.7 & 1.0507 & 165.5 \\
        Teng--Yao & 1.0005 & 19.2 & 1.0034 & 16.5 \\
        Paluch--Elbassioni--van~Zuylen & 1.0000 & 9.0 & 1.0000 & 8.6 \\
        Kosaraju--Park--Stein ($\frac{38}{63}$) & 1.0033 & 249.2 & 1.0248 & 50.2 \\
        Breslauer--Jiang--Jiang $2\frac{2}{3}$ & 1.0007 & 11.5 & 1.0123 & 6.6 \\
        Breslauer--Jiang--Jiang $2\frac{25}{42}$ & 1.0007 & 11.5 & 1.0122 & 6.6 \\
        \bottomrule
    \end{tabular}
\end{table}


\section{Conclusions}

Read on the small instances alone, the comparison says that the choice of
algorithm hardly matters. Across the four experiments of
Sections~\ref{sec:bench_exact} and earlier, every method except the two
Breslauer--Jiang--Jiang variants stays within roughly two percent of the
best result found, and the ordering by measured quality bears no relation
to the ordering by provable constant. Paluch--Elbassioni--van~Zuylen leads
at 1.0018, but \textsc{Greedy} reaches 1.0052 with no
approximation machinery at all, and the gap between them costs a few tens
of milliseconds either way.

Section~\ref{sec:bench_large} shows that this reading is an artefact of
scale. On $n = 1600$ fragments of natural text the two algorithms
still return effectively the same string, differing by 1
character, but \textsc{Greedy} now needs 69 minutes against
9.0 seconds for Paluch--Elbassioni--van~Zuylen, a factor of
463. The empirical exponents behind that gap are 3.0
and roughly 2.0, so it widens with every doubling of
$n$. Which algorithm to use is therefore settled on running time rather
than on length, and it is a difference the small instances cannot show,
because there the same gap is a few tens of milliseconds.

The same correction applies to the algorithms of
Chapter~\ref{chap:breslauer}, which look worst in the chapter and are not.
They average 1.0772 on the small instances and 1.3158 in the
worst case measured, far behind everything else, because their Step~3
chooses cycle representatives $t_C$ so that the subsequent bound can be
proved, and that choice gives up overlap a greedy merge would have taken.
The overlap given up is fixed per cycle, so it dominates when there are few
cycles and amortises when there are many. At $n = 1600$ the same two
algorithms come in at 1.0007, within a thousandth of the best
result, in under twelve seconds. Their weakness was a property of the
instance sizes, not of the construction.

Inside Kosaraju--Park--Stein the picture is different and does not improve
with scale. Section~\ref{sec:bench_kosaraju_versions} shows that Version~1,
the one-line rule discarding the lightest edge of every cycle, is
0.06\,\% off the full algorithm while running 13
times faster, so the path-coloring construction of
Chapter~\ref{chap:path_coloring} and the piece-recombination of Version~3
contribute almost nothing to the measured output. These two versions exist
to hold up the $\frac{38}{63}$ bound, not to shorten strings.

Was implementing these algorithms worth it? On the evidence collected here,
yes, though not for the reason the approximation ratios suggest. None of
them returns a noticeably shorter superstring than a maximum-overlap greedy
merge, and on small inputs none of them needs to. What they provide is a
construction that stays affordable when the input grows, which is exactly
where the greedy heuristic stops being usable, and a guarantee that holds on
inputs no benchmark can be asked to generate. The measurements add the part
the proofs leave open, namely how much of the worst-case pessimism survives
contact with realistic data. Very little of it does, and the practical
argument for these algorithms turns out to rest on their running time rather
than on the bounds they were designed to establish.
